% Options for packages loaded elsewhere
% Options for packages loaded elsewhere
\PassOptionsToPackage{unicode}{hyperref}
\PassOptionsToPackage{hyphens}{url}
\PassOptionsToPackage{dvipsnames,svgnames,x11names}{xcolor}
%
\documentclass[
  letterpaper,
  DIV=11,
  numbers=noendperiod]{scrreprt}
\usepackage{xcolor}
\usepackage{amsmath,amssymb}
\setcounter{secnumdepth}{5}
\usepackage{iftex}
\ifPDFTeX
  \usepackage[T1]{fontenc}
  \usepackage[utf8]{inputenc}
  \usepackage{textcomp} % provide euro and other symbols
\else % if luatex or xetex
  \usepackage{unicode-math} % this also loads fontspec
  \defaultfontfeatures{Scale=MatchLowercase}
  \defaultfontfeatures[\rmfamily]{Ligatures=TeX,Scale=1}
\fi
\usepackage{lmodern}
\ifPDFTeX\else
  % xetex/luatex font selection
\fi
% Use upquote if available, for straight quotes in verbatim environments
\IfFileExists{upquote.sty}{\usepackage{upquote}}{}
\IfFileExists{microtype.sty}{% use microtype if available
  \usepackage[]{microtype}
  \UseMicrotypeSet[protrusion]{basicmath} % disable protrusion for tt fonts
}{}
\makeatletter
\@ifundefined{KOMAClassName}{% if non-KOMA class
  \IfFileExists{parskip.sty}{%
    \usepackage{parskip}
  }{% else
    \setlength{\parindent}{0pt}
    \setlength{\parskip}{6pt plus 2pt minus 1pt}}
}{% if KOMA class
  \KOMAoptions{parskip=half}}
\makeatother
% Make \paragraph and \subparagraph free-standing
\makeatletter
\ifx\paragraph\undefined\else
  \let\oldparagraph\paragraph
  \renewcommand{\paragraph}{
    \@ifstar
      \xxxParagraphStar
      \xxxParagraphNoStar
  }
  \newcommand{\xxxParagraphStar}[1]{\oldparagraph*{#1}\mbox{}}
  \newcommand{\xxxParagraphNoStar}[1]{\oldparagraph{#1}\mbox{}}
\fi
\ifx\subparagraph\undefined\else
  \let\oldsubparagraph\subparagraph
  \renewcommand{\subparagraph}{
    \@ifstar
      \xxxSubParagraphStar
      \xxxSubParagraphNoStar
  }
  \newcommand{\xxxSubParagraphStar}[1]{\oldsubparagraph*{#1}\mbox{}}
  \newcommand{\xxxSubParagraphNoStar}[1]{\oldsubparagraph{#1}\mbox{}}
\fi
\makeatother


\usepackage{longtable,booktabs,array}
\usepackage{calc} % for calculating minipage widths
% Correct order of tables after \paragraph or \subparagraph
\usepackage{etoolbox}
\makeatletter
\patchcmd\longtable{\par}{\if@noskipsec\mbox{}\fi\par}{}{}
\makeatother
% Allow footnotes in longtable head/foot
\IfFileExists{footnotehyper.sty}{\usepackage{footnotehyper}}{\usepackage{footnote}}
\makesavenoteenv{longtable}
\usepackage{graphicx}
\makeatletter
\newsavebox\pandoc@box
\newcommand*\pandocbounded[1]{% scales image to fit in text height/width
  \sbox\pandoc@box{#1}%
  \Gscale@div\@tempa{\textheight}{\dimexpr\ht\pandoc@box+\dp\pandoc@box\relax}%
  \Gscale@div\@tempb{\linewidth}{\wd\pandoc@box}%
  \ifdim\@tempb\p@<\@tempa\p@\let\@tempa\@tempb\fi% select the smaller of both
  \ifdim\@tempa\p@<\p@\scalebox{\@tempa}{\usebox\pandoc@box}%
  \else\usebox{\pandoc@box}%
  \fi%
}
% Set default figure placement to htbp
\def\fps@figure{htbp}
\makeatother





\setlength{\emergencystretch}{3em} % prevent overfull lines

\providecommand{\tightlist}{%
  \setlength{\itemsep}{0pt}\setlength{\parskip}{0pt}}



 


\KOMAoption{captions}{tableheading}
\makeatletter
\@ifpackageloaded{bookmark}{}{\usepackage{bookmark}}
\makeatother
\makeatletter
\@ifpackageloaded{caption}{}{\usepackage{caption}}
\AtBeginDocument{%
\ifdefined\contentsname
  \renewcommand*\contentsname{Table of contents}
\else
  \newcommand\contentsname{Table of contents}
\fi
\ifdefined\listfigurename
  \renewcommand*\listfigurename{List of Figures}
\else
  \newcommand\listfigurename{List of Figures}
\fi
\ifdefined\listtablename
  \renewcommand*\listtablename{List of Tables}
\else
  \newcommand\listtablename{List of Tables}
\fi
\ifdefined\figurename
  \renewcommand*\figurename{Figure}
\else
  \newcommand\figurename{Figure}
\fi
\ifdefined\tablename
  \renewcommand*\tablename{Table}
\else
  \newcommand\tablename{Table}
\fi
}
\@ifpackageloaded{float}{}{\usepackage{float}}
\floatstyle{ruled}
\@ifundefined{c@chapter}{\newfloat{codelisting}{h}{lop}}{\newfloat{codelisting}{h}{lop}[chapter]}
\floatname{codelisting}{Listing}
\newcommand*\listoflistings{\listof{codelisting}{List of Listings}}
\makeatother
\makeatletter
\makeatother
\makeatletter
\@ifpackageloaded{caption}{}{\usepackage{caption}}
\@ifpackageloaded{subcaption}{}{\usepackage{subcaption}}
\makeatother
\usepackage{bookmark}
\IfFileExists{xurl.sty}{\usepackage{xurl}}{} % add URL line breaks if available
\urlstyle{same}
\hypersetup{
  pdftitle={Andrew Tedjapratama},
  pdfauthor={13523148 Andrew Tedjapratama},
  colorlinks=true,
  linkcolor={blue},
  filecolor={Maroon},
  citecolor={Blue},
  urlcolor={Blue},
  pdfcreator={LaTeX via pandoc}}


\title{Andrew Tedjapratama}
\usepackage{etoolbox}
\makeatletter
\providecommand{\subtitle}[1]{% add subtitle to \maketitle
  \apptocmd{\@title}{\par {\large #1 \par}}{}{}
}
\makeatother
\subtitle{Portfolio Asesmen II-2100 KIPP}
\author{13523148 Andrew Tedjapratama}
\date{2025-10-31}
\begin{document}
\maketitle

\renewcommand*\contentsname{Table of contents}
{
\hypersetup{linkcolor=}
\setcounter{tocdepth}{2}
\tableofcontents
}

\bookmarksetup{startatroot}

\chapter*{Selamat Berjumpa}\label{selamat-berjumpa}
\addcontentsline{toc}{chapter}{Selamat Berjumpa}

\markboth{Selamat Berjumpa}{Selamat Berjumpa}

\begin{figure}[H]

{\centering \includegraphics[width=0.95\linewidth,height=\textheight,keepaspectratio]{images/ATDJ.png}

}

\caption{About Me}

\end{figure}%

Hai semua! 👋

Andrew Tedjapratama adalah seorang mahasiswa Teknik Informatika angkatan
2023 di Sekolah Teknik Elektro dan Informatika ITB. Saya adalah
seseorang yang senang belajar hal baru, bekerja dengan ritme yang
teratur, suka ngopi, dan suka ga tidur. Suatu kepercayaan yang dimiliki
saya yaitu percaya kalau hal kecil yang dilakukan dengan niat baik dan
dengan taat kepada Tuhan, bisa berdampak besar, entah di pekerjaan,
percakapan, atau hubungan antarorang.

Situs ini adalah autobiografi singkat untuk memperkenalkan diri saya dan
juga adalah bagian dari UTS mata kuliah II-2100 Komunikasi Interpersonal
dan Publik (KIPP). Lewat laman ini, aku mencoba memperkenalkan diri
bukan lewat prestasi atau pencapaian, tapi lewat cara berpikir dan cara
melihat dunia. Kadang serius, kadang santai --- tapi selalu berusaha
jujur.

Semoga apa yang tertulis di sini bisa memberi sedikit perspektif baru,
atau setidaknya, bikin anda merasa bahwa kita semua masih sama-sama
belajar.

Versi lengkapnya bisa anda lihat di bagian navigasi di sebelah kiri.

Selamat menjelajahi dan have a good day!

\begin{center}\rule{0.5\linewidth}{0.5pt}\end{center}

\bookmarksetup{startatroot}

\chapter{UTS-1 All About Me}\label{uts-1-all-about-me}

\begin{figure}[H]

{\centering \includegraphics[width=0.65\linewidth,height=\textheight,keepaspectratio]{All_About_me/../images/ATDJ2.png}

}

\caption{About Me}

\end{figure}%

\section{Siapa Saya}\label{siapa-saya}

Kalau ditanya siapa saya, mungkin saya bukan tipe orang yang langsung
tahu jawabannya.

Saya lebih suka berpikir lama, bukan karena bingung, tapi karena saya
tahu, semakin kita mencari tahu, semakin kita sadar bahwa ada lebih
banyak hal yang belum kita tahu. Lucunya, justru dari situlah saya
merasa hidup jadi menarik Kadang terlalu logis sampai orang kira saya
lagi cari ribut, padahal cuma penasaran aja. Saya suka diskusi yang
dalam, tentang pilihan, nilai, bahkan kadang hal sefilosofis ``kenapa
kita peduli dengan opini orang lain?''. Bukan buat menang argumen, tapi
buat memahami.

Buat saya, berpikir itu cara paling jujur untuk menghormati hidup karena
hidup bukan cuma untuk dijalani, tapi juga untuk dimengerti. Saya juga
menjunjung tinggi kejujuran dan ketulusan. Saya lebih respek ke orang
yang jujur meskipun belum sempurna, daripada yang berprestasi tapi
palsu.

Ketulusan itu, menurut saya, bukan soal kata-kata manis, tapi soal
sikap, kalau salah, ngaku, atau kalau gak tahu, bilang gak tahu
dibanding membohongi diri.\\
Karena bagi saya, menjadi manusia yang genuine lebih sulit (dan lebih
berharga) daripada sekadar menjadi ``baik''.

\begin{center}\rule{0.5\linewidth}{0.5pt}\end{center}

\section{Nilai yang Saya Pegang}\label{nilai-yang-saya-pegang}

Di luar itu, saya suka hal-hal yang sederhana: ngopi, denger musik,
jalan tanpa tujuan, kadang nulis hal random yang gak pernah saya
publish, atau pun melakukan journalling. Hal-hal sepele kayak gitu
justru sering nyelametin kewarasan saya.. Saya percaya Tuhan bisa
berbicara lewat kesunyian, lewat hujan sore, lagu yang kebetulan pas,
atau percakapan kecil yang datang di waktu yang tepat. Itu momen yang
bikin saya ingat kalau dunia ini gak sepenuhnya chaos.

Saya pernah ngerasain gagal, berkali-kali.\\
Kadang stres, kadang kehilangan arah, tapi saya selalu punya kalimat
yang menenangkan diri:\\
\textgreater{} \emph{``I know I can get there, some way, somehow.''}\\
Saya percaya selama masih mau berjuang dan taat sama Tuhan, jalannya
selalu ada, bahkan kalau sekarang belum kelihatan.

\begin{center}\rule{0.5\linewidth}{0.5pt}\end{center}

\section{Tentang Teman Saya -
Kegagalan}\label{tentang-teman-saya---kegagalan}

Lucunya, saya malah kadang nyari motivasi dari orang-orang yang
teriak-teriak kayak David Goggins atau Andrew Tate.Gunanya bukan sekedar
untuk motivasi, tetapi untuk memakai dan mengarahkan emosi \& energi
kita untuk hal yang lebih positif, ketika melakukan sesuatu yang sulit
atau setelah mengalami gagal. Bukan karena saya mau jadi mereka atau
mengidolakan mereka, tapi karena saya suka cara mereka ngomong tentang
mental toughness, bahwa hidup itu gak akan jadi lebih mudah, tapi kita
bisa jadi lebih kuat. Saya juga menggunakan doa sebagai obat \#1 untuk
refleksi dan ketenangan.

Banyak kegagalan yang akan kita alami kedepannya, dan makin banyak orang
yang akan tidak kita suka, tetapi saya selalu percaya untuk selalu
meningkatkan mutu diri dan meningkatkan kehormatan diri, tanpa turun ke
level yang lebih rendah dan mencoba memaafkan musuh-musuh saya.

Saya bukan orang paling sabar, paling konsisten, atau paling fokus.
Kadang saya perfeksionis, overthinking, dan kebanyakan mikir hal yang
orang lain gak peduli. Tapi di balik itu semua, saya cuma pengen jadi
versi saya yang lebih baik dari kemarin, bukan buat dipuji, tapi buat
bisa menghormati diri sendiri.

\begin{center}\rule{0.5\linewidth}{0.5pt}\end{center}

\section{Being the Best Version}\label{being-the-best-version}

Kalau ditanya apa yang paling saya pelajari dari semua itu, mungkin ini:

\begin{itemize}
\item
  Hidup gak harus sempurna untuk bermakna.
\item
  Kadang yang penting bukan jalannya lurus, tapi kita terus jalan,
  sambil belajar, sambil jatuh, sambil bersyukur.
\item
  Orang tua selalu menjadi tujuan saya, di sebelah Tuhan, sehingga saya
  akan selalu menyadarkan diri bahwa tujuan akhirnya bukan sukses, atau
  uang, tetapi kehangatan dalam melakukan tanggung jawab absolut.
\end{itemize}

Dan mungkin itu cara paling jujur buat mengenal saya --- bukan dari
hal-hal besar, tapi dari cara saya ngadepin hal-hal kecil, tiap hari,
dengan cara saya sendiri.

\bookmarksetup{startatroot}

\chapter{UTS-2 My Song for You}\label{uts-2-my-song-for-you}

\section{Untuk Siapa Lagu Ini
Ditulis}\label{untuk-siapa-lagu-ini-ditulis}

Lagu ini saya persembahkan untuk dua versi saya:\\
\textbf{saya yang sedang berjuang hari ini}, dan \textbf{saya di masa
depan}, yang mungkin sedang kehilangan arah, tapi masih berusaha
berjalan. Dan juga untuk \textbf{Tuhan}, jangkar yang selalu tenang
bahkan saat saya goyah.

``Con Te Partirò'', adalah lagu yang dibuat oleh seorang buta yaitu
Andrea Bocelli. Menurut saya, lagu ini bukan sekadar lagu perpisahan.

Buat saya, lagu ini bicara tentang \textbf{keberanian untuk melangkah,
meski tidak tahu ke mana.}\\
Tentang melepaskan, percaya, dan berjalan bersama Tuhan, bahkan ketika
peta hidup terasa kabur.

\begin{center}\rule{0.5\linewidth}{0.5pt}\end{center}

\section{\texorpdfstring{Lagu: \emph{Con Te Partirò -- Andrea
Bocelli}}{Lagu: Con Te Partirò -- Andrea Bocelli}}\label{lagu-con-te-partiruxf2-andrea-bocelli}

\begin{center}\rule{0.5\linewidth}{0.5pt}\end{center}

\section{Lirik \& Terjemahan (Dibawah)}\label{lirik-terjemahan-dibawah}

\subsection{\texorpdfstring{🇮🇹 \emph{Con Te
Partirò}}{🇮🇹 Con Te Partirò}}\label{con-te-partiruxf2}

Quando sono solo\\
Sogno all'orizzonte\\
E mancan le parole

Sì, lo so che non c'è luce\\
In una stanza quando manca il sole\\
Se non ci sei tu con me, con me

Su le finestre\\
Mostra a tutti il mio cuore\\
Che hai acceso\\
Chiudi dentro me\\
La luce che hai incontrato per strada

Con te partirò\\
Paesi che non ho mai veduto e vissuto con te\\
Adesso sì li vivrò

Con te partirò\\
Su navi per mari\\
Che io lo so\\
No, no, non esistono più\\
Con te io li rivivrò

\begin{center}\rule{0.5\linewidth}{0.5pt}\end{center}

\subsection{Terjemahan Indonesia}\label{terjemahan-indonesia}

Saat aku sendiri,\\
Aku bermimpi tentang cakrawala ---\\
dan kata-kata pun tak lagi ada.

Aku tahu,\\
tak ada cahaya di dalam ruangan tanpa matahari,\\
jika kamu tak ada bersamaku.

Bukalah jendela,\\
tunjukkan pada dunia hati yang telah kamu nyalakan,\\
dan biarkan cahaya itu tinggal dalam diriku.

Bersamamu, aku akan pergi,\\
ke tempat-tempat yang belum pernah kulihat atau kualami.\\
Sekarang aku akan menjalaninya.

Bersamamu, aku akan berlayar,\\
menyebrangi lautan yang dulu tak ada,\\
dan bersamamu, aku akan menghidupkannya kembali.

\begin{center}\rule{0.5\linewidth}{0.5pt}\end{center}

\section{Makna Pribadi}\label{makna-pribadi}

Saya tidak menganggap lagu ini tentang cinta antar manusia. Buat saya,
lagu ini tentang \textbf{iman}. Tentang kepercayaan yang pelan tapi
pasti. Setiap kali Andrea Bocelli menyanyikan \emph{``Con te partirò''},
saya selalu membayangkan Tuhan berkata, ``Kamu tidak sendiri --- Aku
ikut berjalan bersamamu.''

Nada lagunya megah, tapi tidak sombong;\\
seakan mengingatkan saya bahwa \emph{percaya} bukan berarti tahu
segalanya,\\
melainkan \textbf{berani melangkah walau belum tahu hasilnya}.

\begin{center}\rule{0.5\linewidth}{0.5pt}\end{center}

\section{Refleksi}\label{refleksi}

Kalau suatu saat saya membaca halaman ini lagi ---\\
mungkin saya sudah berbeda: lebih tua, lebih rasional, mungkin juga
lebih keras kepala.\\
Tapi saya harap saya tidak lupa makna dari lagu ini.

Bahwa setiap perjalanan punya waktunya sendiri.\\
Bahwa saya pernah percaya, di tengah bingung dan takut,\\
bahwa ada Tuhan yang tetap berjalan di samping saya.

\begin{itemize}
\tightlist
\item
  ``Con Te Partirò'' bukan sekadar \emph{``bersamamu aku akan
  pergi.''}\\
\item
  Tapi juga \emph{``aku akan percaya, dan tidak takut untuk
  melangkah.''}
\end{itemize}

Dan buat saya, itu doa yang sederhana tapi paling jujur.\\
Doa dari seseorang yang gak punya semua jawabannya,\\
tapi masih berani berjalan, dengan iman yang kecil, tapi tulus.

\bookmarksetup{startatroot}

\chapter{UTS-3 My Stories for You}\label{uts-3-my-stories-for-you}

\section{Untuk Mama \& Papa}\label{untuk-mama-papa}

Papa dan mama selalu jadi dua sosok yang paling saya kagumi tapi juga
paling sering saya diamkan. Dulu saya kira kasih sayang itu harus selalu
terlihat dari kata-kata, tapi makin ke sini saya sadar, cinta orang tua
sering hadir dalam bentuk yang tidak berbunyi, dari cara mama nanyain
makan tanpa terlihat bawel, atau papa yang diam-diam memastikan semuanya
berjalan baik. Kalian gak pernah minta dibalas, tapi tanpa sadar semua
yang saya lakukan hari ini selalu ada jejak kalian di belakangnya.

Maaf kalau saya terlalu naif dan egois selama ini, mungkin hidup adalah
menyadari betapa perlunya kita bersyukur. Mau itu kepada mama papa
maupun kepada Tuhan, karena saya percaya mama papa adalah perantara
Tuhan. Saya minta maaf atas kata kasar dan kekurangan saya. Saya tahu
saya gak selalu ekspresif, kadang malah terlalu sibuk dengan hal-hal
kecil yang bahkan gak penting. Tapi di antara semua kekacauan itu,
kalian selalu jadi alasan kenapa saya pengen jadi lebih baik. Kalian
yang bikin saya belajar tentang tanggung jawab, disiplin, dan keikhlasan
tanpa pamrih. Kalau suatu hari nanti saya bisa mencapai sesuatu, itu
bukan karena saya luar biasa, tapi karena saya dibentuk oleh dua orang
yang gak pernah berhenti percaya pada saya.

\section{Untuk Koko}\label{untuk-koko}

Kita jarang ketemu, karena koko ada di luar negeri sejak SMP, umur kita
jauh, dan saya tahu kita beda banget dalam cara berpikir. Tapi entah
kenapa, perbedaan itu justru bikin saya belajar banyak hal. Kamu mungkin
gak sadar, tapi cara kamu menghadapi dunia dengan tenang itu sering jadi
cermin buat saya. Kadang saya yang terlalu keras kepala, kadang kamu
yang terlalu santai, tapi di tengah semua itu, saya selalu tahu --- kamu
adalah orang yang saya hormati, bahkan kalau saya gak selalu bilang.
Saya selalu melihat Koko sebagai orang yang paling sabar, positif, dan
baik.

Saya juga belajar dari cara koko nunjukin perhatian tanpa ngomong
banyak. Dari tindakan kecil yang kelihatannya biasa aja, tapi berarti
besar. Kamu ngajarin saya bahwa peduli gak selalu harus lewat omongan
panjang, tapi lewat kehadiran yang tulus. Terima kasih udah tetap jadi
``koko'' yang apa adanya, bukan yang sempurna, tapi yang nyata. Dan saya
rasa, itu lebih dari cukup. Looking forward to our lives next.

\section{Untuk Teman-Teman Dekat Saya (teman lama, teman sekarang, teman
di masa
depan)}\label{untuk-teman-teman-dekat-saya-teman-lama-teman-sekarang-teman-di-masa-depan}

Kalian mungkin gak tahu seberapa besar pengaruh kalian dalam perjalanan
saya. Saya bukan tipe orang yang gampang terbuka, tapi di antara tawa,
obrolan receh, dan malam-malam absurd, saya merasa diterima apa adanya.
Kadang kita bahas hal seaneh teori konspirasi atau se-serius masa depan,
tapi semuanya terasa ringan. Bersama kalian, saya bisa berhenti mikir
terlalu jauh dan belajar menikmati detik yang ada di depan mata.

Kalian juga orang-orang yang ngingetin saya buat gak terlalu keras sama
diri sendiri. Bahwa gagal itu gak selalu berarti berhenti, dan
kehilangan arah bukan dosa. Kalian datang di momen yang tepat --- waktu
saya mulai lupa gimana rasanya hidup tanpa beban. Terima kasih udah
tetap ada, bahkan ketika saya gak selalu hadir. Kalian bagian dari versi
terbaik saya yang sedang saya bentuk pelan-pelan.

\section{Untuk Guru-Guru Saya}\label{untuk-guru-guru-saya}

Saya dulu pikir guru itu cuma orang yang ngasih pelajaran, nilai, dan
tugas. Tapi makin dewasa, saya sadar guru itu jauh lebih dari itu ---
mereka adalah orang yang menanamkan cara berpikir. Dari setiap koreksi,
teguran, dan nasihat, saya pelan-pelan belajar bahwa hidup bukan cuma
tentang benar atau salah, tapi tentang berproses dengan rendah hati.
Saya sering kali lambat paham, tapi kalian tetap sabar membimbing tanpa
menghakimi.

Saya juga mau berterima kasih karena banyak dari kalian yang ngajarin
lewat contoh, bukan cuma teori. Kalian menunjukkan bahwa jadi manusia
yang bijak bukan berarti tanpa cacat, tapi tetap berdiri tegak di tengah
kekurangan. Dari kalian, saya belajar bahwa integritas dan empati gak
bisa diajarkan lewat kata-kata, tapi lewat keteladanan. Dan saya harap,
satu hari nanti, saya bisa jadi seseorang yang memberikan arti dengan
cara yang sama.

\bookmarksetup{startatroot}

\chapter{UTS-4 My SHAPE (Spiritual Gifts, Heart, Abilities, Personality,
Experiences)}\label{uts-4-my-shape-spiritual-gifts-heart-abilities-personality-experiences}

\begin{center}\rule{0.5\linewidth}{0.5pt}\end{center}

\section{0) Ringkasan 1 Halaman}\label{ringkasan-1-halaman}

\textbf{Peran Inti:} Mahasiswa, pemikir reflektif, dan pembelajar seumur
hidup yang terus berusaha memahami makna di balik setiap hal.\\
\textbf{Misi:} Menggabungkan logika, kejujuran, dan iman untuk
menciptakan dampak positif , baik dalam kerja, relasi, maupun
pelayanan.\\
\textbf{Kekuatan Utama:} rasa ingin tahu yang tinggi, kemampuan berpikir
logis-sistematis, empati reflektif, serta kemauan untuk terus
memperbaiki diri.\\
\textbf{Dampak yang Dituju:} menjadi pribadi yang bisa membawa
ketenangan, kejelasan, dan semangat bertumbuh bagi orang di sekeliling
saya.

\textbf{Peta SHAPE (singkat):} - \textbf{S , Spiritual Gifts:} Teaching,
Discernment, Encouragement, Leadership.\\
- \textbf{H , Heart:} nilai kejujuran, kedisiplinan, refleksi, dan
kesetiaan dalam iman.\\
- \textbf{A , Abilities:} analisis logika, pemrograman, komunikasi,
problem solving, dan menulis ide secara runtut.\\
- \textbf{P , Personality:} reflektif-strategis, tenang, perfeksionis,
dan suka berpikir mendalam.\\
- \textbf{E , Experiences:} melalui kegagalan, stres, pencarian iman,
dan perjalanan akademik yang membentuk ketahanan mental.

\begin{center}\rule{0.5\linewidth}{0.5pt}\end{center}

\section{1) S - Spiritual Gifts (Karunia
Rohani)}\label{s---spiritual-gifts-karunia-rohani}

Saya menyadari bahwa kemampuan mengajar dan berpikir logis bukan hanya
bakat, tetapi juga karunia untuk membantu orang memahami sesuatu dengan
jernih.\\
Saya sering menemukan kedamaian dalam diskusi , terutama saat bisa
membuat orang lain memahami hal rumit dengan sederhana.\\
Selain itu, saya punya dorongan kuat untuk meneguhkan orang lain,
terutama teman-teman yang sedang kehilangan arah.\\
Saya melihat kepemimpinan bukan soal memerintah, tetapi soal memberi
arah dan menjaga ritme agar tim tetap berjalan bersama.

\textbf{Indikator Bukti:} mentoring teman dalam proyek koding, sharing
reflektif di diskusi kampus, dan bimbingan kelompok belajar.

\begin{center}\rule{0.5\linewidth}{0.5pt}\end{center}

\section{2) H - Heart (Minat, Nilai,
Kepedulian)}\label{h---heart-minat-nilai-kepedulian}

Saya peduli pada kejujuran, keaslian, dan pencarian makna hidup yang
nyata. Saya percaya setiap hal , bahkan yang gagal , bisa jadi bahan
belajar jika dilihat dengan hati yang terbuka.\\
Hal yang paling saya cintai adalah proses berpikir dan berbagi
pandangan; bukan untuk terlihat pintar, tapi untuk menemukan kebenaran
bersama.\\
Saya juga punya ketertarikan besar pada musik, tulisan reflektif, dan
teknologi yang bisa digunakan untuk memperbaiki kehidupan manusia.

\textbf{Masalah yang ingin dipecahkan:} bagaimana cara manusia tetap
tulus dan sadar nilai di tengah dunia yang serba cepat dan kompetitif.

\begin{center}\rule{0.5\linewidth}{0.5pt}\end{center}

\section{3) A - Abilities (Kemampuan
Andal)}\label{a---abilities-kemampuan-andal}

\begin{itemize}
\tightlist
\item
  \textbf{Hard Skills:} Programming (JavaScript, React, Python), UI/UX,
  data analysis, logical problem-solving, technical writing.\\
\item
  \textbf{Soft Skills:} komunikasi jujur, mendengarkan aktif, disiplin,
  manajemen stres, dan kemampuan adaptasi.\\
  Saya bukan orang yang paling cepat belajar, tapi saya konsisten. Saya
  melihat kemampuan berpikir sistematis sebagai alat untuk membawa
  keteraturan pada kekacauan.
\end{itemize}

\textbf{Aplikasi nyata:} pembuatan proyek web interaktif pribadi dan
inisiatif membantu teman memahami konsep coding dengan cara sederhana.

\begin{center}\rule{0.5\linewidth}{0.5pt}\end{center}

\section{4) P - Personality (Gaya Kerja \&
Kolaborasi)}\label{p---personality-gaya-kerja-kolaborasi}

Saya tergolong \textbf{INTJ-T}, tipe yang reflektif, analitis, dan
strategis. Saya cenderung fokus pada gambaran besar namun tetap
detail-oriented.\\
Saya bukan tipe yang suka spotlight, tapi saya berusaha menjaga standar
tinggi dalam setiap hal yang saya kerjakan.\\
Dalam tim, saya sering jadi orang yang menenangkan, terutama ketika
suasana mulai chaos.\\
Saya lebih suka hasil yang jujur daripada pencitraan sempurna.

\begin{center}\rule{0.5\linewidth}{0.5pt}\end{center}

\section{5) E - Experiences (Pengalaman
Pembentuk)}\label{e---experiences-pengalaman-pembentuk}

Saya dibentuk lewat banyak stres dan kegagalan , proyek yang gagal,
waktu yang terbuang, dan rasa kehilangan arah.\\
Namun, dari semua itu saya belajar mengenal Tuhan dengan cara yang lebih
pribadi.\\
Saya belajar bahwa keberhasilan tanpa kedewasaan hati tidak berarti
apa-apa.\\
Doa dan refleksi menjadi cara saya untuk tetap waras dan menyadari bahwa
semua proses, sekecil apa pun, punya tujuan.

\textbf{Pelajaran inti:} ketekunan lebih penting daripada kecepatan;
kejujuran lebih penting daripada pencapaian; Tuhan tetap jadi fondasi di
tengah semua perubahan.

\begin{center}\rule{0.5\linewidth}{0.5pt}\end{center}

\section{6) Piagam Diri (Self-Charter)}\label{piagam-diri-self-charter}

\textbf{Misi Hidup:} menjadi pribadi yang menggabungkan rasionalitas dan
iman untuk menghadirkan kejujuran, kedamaian, dan pertumbuhan di
lingkungan saya.\\
\textbf{Nilai Inti:} integritas, keaslian, rasa hormat, refleksi,
ketekunan.\\
\textbf{Peran Inti:} pemikir yang membangun, bukan meruntuhkan;
pembelajar yang menginspirasi lewat contoh; rekan yang bisa dipercaya.\\
\textbf{Kompas Keputusan:}\\
1. Apakah ini membawa kebaikan bagi orang lain?\\
2. Apakah ini selaras dengan hati nurani dan iman saya?\\
3. Apakah ini bisa saya pertanggungjawabkan dengan jujur?\\
\textbf{Janji Pelayanan:} tetap hadir dengan rendah hati, mendengarkan
sebelum menilai, dan berusaha memperbaiki sistem, bukan orangnya.\\
\textbf{Batasan:} menolak hal yang melanggar etika, menjauh dari
lingkungan toksik, dan menjaga ritme antara kerja, iman, dan keluarga.

\begin{center}\rule{0.5\linewidth}{0.5pt}\end{center}

\section{7) Narasi 90 Detik (Elevator
Pitch)}\label{narasi-90-detik-elevator-pitch}

Saya Andrew Tedjapratama, seorang mahasiswa yang percaya bahwa berpikir
dan beriman tidak harus bertentangan.\\
Saya ingin menggunakan logika dan empati untuk membantu orang lain
melihat hidup dengan lebih jernih.\\
Kekuatan saya adalah refleksi, disiplin, dan kemampuan mengolah stres
menjadi pembelajaran.\\
Saya percaya Tuhan bekerja lewat proses, dan saya ingin menjalani setiap
langkah hidup dengan niat baik dan tanggung jawab , bukan untuk terlihat
hebat, tapi untuk tetap jujur dan berguna.

\begin{center}\rule{0.5\linewidth}{0.5pt}\end{center}

\section{8) Service-Fit Map (Tempat Saya Paling
Berdampak)}\label{service-fit-map-tempat-saya-paling-berdampak}

\begin{itemize}
\tightlist
\item
  \textbf{Kampus:} diskusi reflektif, proyek teknologi kolaboratif,
  mentorship rekan.\\
\item
  \textbf{Komunitas:} kegiatan pelayanan dan refleksi iman.\\
\item
  \textbf{Kreativitas:} menulis narasi reflektif, desain web, atau media
  visual bermakna.\\
\item
  \textbf{Pribadi:} ruang sunyi untuk berpikir, berdoa, dan memahami
  diri agar bisa memberi lebih banyak pada orang lain.
\end{itemize}

\begin{center}\rule{0.5\linewidth}{0.5pt}\end{center}

\section{9) Evidences (Artefak \&
Tautan)}\label{evidences-artefak-tautan}

\begin{itemize}
\tightlist
\item[$\square$]
  Proyek React Portfolio -- GitHub\\
\item[$\square$]
  Tulisan reflektif pribadi\\
\item[$\square$]
  Kontribusi pada proyek kelompok\\
\item[$\square$]
  Catatan jurnal pribadi tentang pengembangan diri
\end{itemize}

\begin{center}\rule{0.5\linewidth}{0.5pt}\end{center}

\section{10) Rencana Aksi 90 Hari
(SMART)}\label{rencana-aksi-90-hari-smart}

\begin{enumerate}
\def\labelenumi{\arabic{enumi}.}
\tightlist
\item
  \textbf{Menyelesaikan semua tugas UTS dan UAS tepat waktu.}\\
  \emph{Outcome:} semua artefak KIPP terpublikasi; \emph{Due:} T-14.\\
\item
  \textbf{Membangun proyek mini-web pribadi.}\\
  \emph{Outcome:} 1 laman aktif; \emph{Due:} T-45.\\
\item
  \textbf{Menulis 3 refleksi pribadi tentang nilai \& iman.}\\
  \emph{Outcome:} 3 tulisan selesai; \emph{Due:} T-75.\\
\item
  \textbf{Membiasakan journaling malam untuk refleksi diri.}\\
  \emph{Outcome:} kebiasaan konsisten; \emph{Due:} T-90.
\end{enumerate}

\begin{center}\rule{0.5\linewidth}{0.5pt}\end{center}

\section{11) SHAPE - CPMK (Keterkaitan
Pembelajaran)}\label{shape---cpmk-keterkaitan-pembelajaran}

\begin{itemize}
\tightlist
\item
  \textbf{Self-awareness (CPMK-1):} tercermin dalam Piagam Diri \&
  refleksi.\\
\item
  \textbf{Empati \& Komunikasi Etis (CPMK-2):} melalui mentoring dan
  interaksi tim.\\
\item
  \textbf{Storytelling (CPMK-3):} dari tulisan reflektif dan narasi
  personal.\\
\item
  \textbf{Kepemimpinan (CPMK-4):} dari cara saya mengambil tanggung
  jawab dan memberi arah di proyek.
\end{itemize}

\begin{center}\rule{0.5\linewidth}{0.5pt}\end{center}

\section{12) Self-Assessment Rubrik
UTS-4}\label{self-assessment-rubrik-uts-4}

\begin{longtable}[]{@{}
  >{\raggedright\arraybackslash}p{(\linewidth - 6\tabcolsep) * \real{0.2308}}
  >{\raggedright\arraybackslash}p{(\linewidth - 6\tabcolsep) * \real{0.2308}}
  >{\centering\arraybackslash}p{(\linewidth - 6\tabcolsep) * \real{0.3077}}
  >{\raggedright\arraybackslash}p{(\linewidth - 6\tabcolsep) * \real{0.2308}}@{}}
\toprule\noalign{}
\begin{minipage}[b]{\linewidth}\raggedright
Kriteria
\end{minipage} & \begin{minipage}[b]{\linewidth}\raggedright
Deskripsi
\end{minipage} & \begin{minipage}[b]{\linewidth}\centering
Skor (1--5)
\end{minipage} & \begin{minipage}[b]{\linewidth}\raggedright
Bukti
\end{minipage} \\
\midrule\noalign{}
\endhead
\bottomrule\noalign{}
\endlastfoot
Kelengkapan SHAPE & S-H-A-P-E jelas \& terisi & \textbf{5} & Ringkasan
0--5 \\
Koherensi Piagam Diri & misi-nilai-peran konsisten & \textbf{5} & Bagian
6 \\
Narasi 90 Detik & ringkas, reflektif, mengundang aksi & \textbf{5} &
Bagian 7 \\
Evidence \& Aksi 90 Hari & rencana SMART \& bukti relevan & \textbf{5} &
Bagian 9--10 \\
\end{longtable}

\textbf{Total:} 20 / 20\\
\textbf{Tingkat:} A (\textgreater= 85 \%)

\begin{center}\rule{0.5\linewidth}{0.5pt}\end{center}

\section{13) Versi Ultra-Ringkas (\textless= 140
kata)}\label{versi-ultra-ringkas-140-kata}

Saya Andrew Tedjapratama, seseorang yang berusaha hidup jujur, berpikir
dalam, dan tetap beriman di tengah kompleksitas hidup.\\
Kekuatan saya ada pada refleksi, logika, dan ketenangan.\\
Saya percaya setiap kegagalan adalah bahan bakar untuk pertumbuhan, dan
setiap proses adalah bentuk kasih Tuhan yang tersembunyi.\\
Misi saya sederhana: terus belajar, bekerja dengan hati, dan menjaga
integritas dalam hal sekecil apa pun.\\
Saya ingin menjadi pribadi yang berguna, bukan besar , karena yang kecil
tapi tulus, sering kali lebih berdampak.

\bookmarksetup{startatroot}

\chapter{UTS-5 My Personal Reviews}\label{uts-5-my-personal-reviews}

\section{6.1 Identifikasi}\label{identifikasi}

\textbf{Nama Mahasiswa dan NIM:} Andrew Tedjapratama (13523148)\\
\textbf{Nama Penilai:} Self-Assessment

\begin{center}\rule{0.5\linewidth}{0.5pt}\end{center}

\section{6.2 Tinjauan Umum}\label{tinjauan-umum}

Seluruh rangkaian UTS-1 sampai UTS-4 bagi saya adalah refleksi bertahap
tentang bagaimana berpikir dan berkomunikasi secara lebih jujur dan
bermakna.\\
UTS-1 menunjukkan pemahaman diri dan nilai personal, UTS-2 menggambarkan
sisi emosional dan spiritualitas saya, UTS-3 memperlihatkan empati
interpersonal, dan UTS-4 membentuk disiplin berpikir terstruktur.\\
Dari keempatnya, saya melihat konsistensi: setiap karya berangkat dari
kejujuran, berpijak pada refleksi, dan tetap menjaga logika berpikir
yang rapi.

\begin{center}\rule{0.5\linewidth}{0.5pt}\end{center}

\section{6.3 Tinjauan Spesifik}\label{tinjauan-spesifik}

\subsection{\texorpdfstring{6.3.1 UTS-1 --- \emph{All About
Me}}{6.3.1 UTS-1 --- All About Me}}\label{uts-1-all-about-me-1}

\textbf{Konten:} Esai reflektif dengan narasi personal yang logis,
tenang, dan kuat secara nilai. Menggabungkan rasionalitas dan
spiritualitas tanpa terkesan menggurui.\\
\textbf{Penilaian (Rubrik Tabel 2):}\\
- Orisinalitas : 5\\
- Keterlibatan : 5\\
- Humor : 3\\
- Wawasan : 5\\
\textbf{Skor rata-rata:} (5 + 5 + 3 + 5) / 4 = \textbf{4.5 / 5.00}\\
\textbf{Saran:} Sudah solid dan berkarakter. Bisa diperkuat dengan
sedikit variasi tone agar lebih dinamis, tapi arah refleksi sudah
matang.

\begin{center}\rule{0.5\linewidth}{0.5pt}\end{center}

\subsection{\texorpdfstring{6.3.2 UTS-2 --- \emph{My Song for
You}}{6.3.2 UTS-2 --- My Song for You}}\label{uts-2-my-song-for-you-1}

\textbf{Konten:} Lagu \emph{Con Te Partirò} (Andrea Bocelli) dijadikan
refleksi spiritual dan eksistensial tentang perpisahan dan keteguhan
iman. Refleksinya puitis, bermakna, dan menunjukkan kedewasaan
berpikir.\\
\textbf{Penilaian (Rubrik Tabel 3):}\\
- Orisinalitas : 5\\
- Keterlibatan : 5\\
- Humor : 1\\
- Inspirasi : 5\\
\textbf{Skor rata-rata:} (5 + 5 + 1 + 5) / 4 = \textbf{4.0 / 5.00}\\
\textbf{Saran:} Pemilihan lagu dan interpretasi sangat kuat. Tambahkan
sedikit konteks pribadi agar hubungan antara makna lagu dan pengalaman
hidup lebih terasa natural.

\begin{center}\rule{0.5\linewidth}{0.5pt}\end{center}

\subsection{\texorpdfstring{6.3.3 UTS-3 --- \emph{My Stories for
You}}{6.3.3 UTS-3 --- My Stories for You}}\label{uts-3-my-stories-for-you-1}

\textbf{Konten:} Narasi reflektif kepada keluarga, teman, dan guru yang
menggambarkan empati dan rasa terima kasih dengan bahasa yang jujur,
tenang, dan tidak berlebihan.\\
\textbf{Penilaian (Rubrik Tabel 4):}\\
- Orisinalitas : 5\\
- Keterlibatan : 5\\
- Pengembangan Narasi : 5\\
- Inspirasi : 5\\
\textbf{Skor rata-rata:} (5 + 5 + 5 + 5) / 4 = \textbf{5.0 / 5.00}\\
\textbf{Saran:} Karya paling kuat dan paling otentik secara emosional.
Gaya narasinya alami dan sangat efektif menyampaikan rasa hormat tanpa
melodrama.

\begin{center}\rule{0.5\linewidth}{0.5pt}\end{center}

\subsection{\texorpdfstring{6.3.4 UTS-4 --- \emph{My
SHAPE}}{6.3.4 UTS-4 --- My SHAPE}}\label{uts-4-my-shape}

\textbf{Konten:} Piagam diri terstruktur yang merangkum misi hidup,
kekuatan utama, nilai inti, dan pengalaman pribadi. Lengkap dan
realistis, memperlihatkan kejelasan arah hidup serta konsistensi
nilai.\\
\textbf{Penilaian (Rubrik Tabel 5):}\\
- Orisinalitas : 5\\
- Keterlibatan : 4\\
- Pengembangan Narasi : 4\\
- Inspirasi : 5\\
\textbf{Skor rata-rata:} (5 + 4 + 4 + 5) / 4 = \textbf{4.5 / 5.00}\\
\textbf{Saran:} Struktur sudah sangat rapi. Hanya perlu sedikit tambahan
narasi transisi agar tiap elemen SHAPE terasa lebih mengalir seperti
cerita.

\begin{center}\rule{0.5\linewidth}{0.5pt}\end{center}

\subsection{\texorpdfstring{6.3.5 UTS-5 --- \emph{My Personal
Review}}{6.3.5 UTS-5 --- My Personal Review}}\label{uts-5-my-personal-review}

\textbf{Konten:} Self-assessment reflektif yang jujur dan argumentatif
dengan analisis berbasis rubrik. Menunjukkan kemampuan evaluasi diri
yang matang.\\
\textbf{Penilaian (Rubrik Tabel 6):}\\
- Pemahaman Konsep : 5\\
- Analisis Kritis Pesan : 5\\
- Argumentasi (Logos) : 5\\
- Etos \& Empati : 5\\
- Rekomendasi Perbaikan : 4\\
\textbf{Skor rata-rata:} (5 + 5 + 5 + 5 + 4) / 5 = \textbf{4.8 / 5.00}\\
\textbf{Saran:} Analisis sudah sangat objektif dan matang. Bisa
diperkuat dengan refleksi ringkas tentang perubahan karakter selama
proses UTS 1--4.

\begin{center}\rule{0.5\linewidth}{0.5pt}\end{center}

\section{6.4 Perhitungan Skor CPMK
(Ringkasan)}\label{perhitungan-skor-cpmk-ringkasan}

\begin{longtable}[]{@{}lccc@{}}
\toprule\noalign{}
UTS & Skor (dari 5) & CPMK-2 (Bobot) & Kontribusi \\
\midrule\noalign{}
\endhead
\bottomrule\noalign{}
\endlastfoot
UTS-1 & 4.50 & 6 & (4.50 / 5) × 6 = 5.40 \\
UTS-2 & 4.00 & 7 & (4.00 / 5) × 7 = 5.60 \\
UTS-3 & 5.00 & 7 & (5.00 / 5) × 7 = 7.00 \\
UTS-4 & 4.50 & 6 & (4.50 / 5) × 6 = 5.40 \\
UTS-5 & 4.80 & 10 & (4.80 / 5) × 10 = 9.60 \\
\textbf{Total} & & & \textbf{33.00 / 36 = 91.7\% (A)} \\
\end{longtable}

\textbf{Download file skor (untuk bukti \& laporan):}\\
Download CSV Ringkasan\\
Download Excel Ringkasan ---

\section{6.5 Kesimpulan Singkat}\label{kesimpulan-singkat}

UTS-1 hingga UTS-5 bukan sekadar dokumentasi perjalanan akademik, tapi
catatan bagaimana saya belajar mengenal diri sendiri.\\
UTS-1 membuka refleksi tentang makna hidup dan nilai; UTS-2 menguatkan
sisi spiritual; UTS-3 menegaskan empati dan hubungan personal; UTS-4
menyatukan semuanya ke dalam struktur yang logis dan terarah.\\
UTS-5 akhirnya menjadi cermin dari semua proses itu --- bukan untuk
mencari nilai, tapi memastikan bahwa saya berkembang sebagai pribadi
yang berpikir kritis, berempati, dan tetap punya arah hidup yang jelas.

\begin{center}\rule{0.5\linewidth}{0.5pt}\end{center}

\textbf{Final Score: 91.7\% (A)} --- \emph{Reflektif, konsisten, dan
memenuhi rubrik dengan kedalaman dan keseimbangan.}

\section{Refleksi}\label{refleksi-1}

Melalui proses ini saya belajar melihat tugas saya sendiri secara lebih
objektif, dan menghargai proses orang lain.\\
Menilai teman membuat saya sadar bahwa setiap karya punya cara unik
dalam mengekspresikan nilai dan perjalanan diri.\\
Sedangkan \emph{self-assessment} membantu saya menyadari kekuatan dan
area yang masih bisa saya tingkatkan, terutama dalam konsistensi dan
kedalaman refleksi. Saya percaya tujuan akhirnya bukan angka, tapi
kejujuran dalam mengenal diri sendiri dan memberi nilai dengan empati.

\bookmarksetup{startatroot}

\chapter{UAS-1 My Concepts}\label{uas-1-my-concepts}

\section{Arsitektur Distribusi Kesehatan: Solusi Sistem Informasi untuk
4.5 Miliar
Penduduk}\label{arsitektur-distribusi-kesehatan-solusi-sistem-informasi-untuk-4.5-miliar-penduduk}

\subsection{Pendahuluan: Fokus pada Masalah Kesehatan Global (Topik
No.~7)}\label{pendahuluan-fokus-pada-masalah-kesehatan-global-topik-no.-7}

Konsep ini dirancang sebagai respons terhadap fakta bahwa saat ini
sekitar 50\% penduduk bumi (4.5 miliar orang) belum mendapatkan layanan
kesehatan esensial secara penuh. Masalah kesehatan global bukan hanya
soal penemuan obat baru, melainkan tantangan logistik dan sistem
informasi agar akses menjadi merata. Dengan memanfaatkan teknologi di
era AI, sistem ini bertujuan untuk memindahkan beban deteksi dini dari
rumah sakit besar ke tingkat komunitas paling dasar.

\begin{center}\rule{0.5\linewidth}{0.5pt}\end{center}

\subsection{Diagnostik Terdistribusi via
Edge-AI}\label{diagnostik-terdistribusi-via-edge-ai}

Masalah utama di banyak negara berkembang adalah rasio tenaga medis yang
sangat timpang dan fasilitas yang terbatas. Menurut saya, kita tidak
bisa hanya mengandalkan pembangunan rumah sakit fisik yang memakan waktu
lama.

\begin{itemize}
\item
  Implementasi Teknis: Menggunakan model TinyML yang ringan untuk
  dijalankan pada smartphone spesifikasi rendah di daerah dengan
  konektivitas terbatas.
\item
  Fungsi Utama: Melakukan triage atau pemilahan gejala awal secara
  otomatis untuk penyakit yang seharusnya bisa dicegah atau diobati.
\item
  Trade-off: Yang menarik di sini yaitu adalah dilema antara akurasi dan
  portabilitas. Trade-off yang saya ambil adalah mengoptimalkan model
  agar bisa berjalan 100\% offline di hardware murah, meskipun harus
  sedikit mengorbankan kedalaman neural network jika dibandingkan dengan
  model yang berjalan di cloud.
\end{itemize}

\subsection{Dashboard Prediktif Logistik Obat
Esensial}\label{dashboard-prediktif-logistik-obat-esensial}

Stagnansi kemajuan Universal Health Coverage (UHC) sejak 2015 sangat
menunjukkan kegagalan dalam menjaga ketersediaan layanan dasar. Banyak
klinik di daerah terpencil mengalami kekosongan stok obat pada saat yang
krusial. - Sistem Dashboard: Membangun integrasi API antara data stok
puskesmas dan gudang logistik pusat secara transparan.

\begin{itemize}
\item
  Analisis Data: Sistem memprediksi lonjakan kebutuhan obat berdasarkan
  tren penyakit musiman, seperti diare yang mengakibatkan kematian 297
  ribu balita setiap tahunnya.
\item
  Personal Voice: Problem di sini sering kali bukan karena stok obat
  tidak ada, tapi karena distribusi yang buta terhadap data lapangan.
  Dari sisi implementasi, ini adalah masalah sinkronisasi data yang
  harusnya bisa diselesaikan dengan sistem informasi yang lebih
  responsif.
\end{itemize}

\subsection{Mitigasi Risiko Finansial bagi Pasien
Rentan}\label{mitigasi-risiko-finansial-bagi-pasien-rentan}

Biaya kesehatan yang tinggi tentunya menjadi faktor utama yang mendorong
jutaan orang jatuh ke jurang kemiskinan setiap tahunnya. Saat ini,
sekitar 2 miliar orang di dunia menghadapi beban finansial berat karena
membayar pengobatan dari kantong sendiri. - Analisis Risiko: Mengingat
minat saya di bidang aktuaria, saya mengusulkan modul penilaian risiko
finansial untuk mengidentifikasi kelompok yang paling rentan terhadap
biaya medis katastrofik.

\begin{itemize}
\item
  Target Dampak: Sistem ini membantu pemerintah untuk melakukan
  intervensi subsidi secara tepat sasaran agar biaya kesehatan tidak
  menghancurkan ekonomi keluarga.
\item
  Realistic View: Saya sadar ini bukan solusi global yang instant juga.
  Implementasinya harus dilakukan bertahap, mulai dari digitalisasi data
  kesehatan primer sebelum melompat ke sistem asuransi mikro yang
  kompleks.
\end{itemize}

\begin{center}\rule{0.5\linewidth}{0.5pt}\end{center}

\subsection{Rencana Implementasi dan
Tahapan}\label{rencana-implementasi-dan-tahapan}

\begin{itemize}
\tightlist
\item
  Fase 1: Digitalisasi data pasien di klinik primer pedesaan untuk
  membangun accuracy data dasar.
\item
  Fase 2: Deploy alat triage AI ringan pada smartphone tenaga medis guna
  mempercepat deteksi dini.
\item
  Fase 3: Integrasi sistem asuransi mikro dan subsidi tepat sasaran
  untuk melindungi pasien dari kemiskinan medis.
\end{itemize}

\begin{center}\rule{0.5\linewidth}{0.5pt}\end{center}

\subsection{Refleksi dan Kesimpulan}\label{refleksi-dan-kesimpulan}

Menurut saya, sebuah sistem informasi baru bisa disebut mahakarya jika
ia mampu berfungsi di tengah keterbatasan lingkungan, bukan hanya di
lingkungan lab yang ideal. AI sering kali dianggap sebagai teknologi
elit, padahal manfaat terjujurnya adalah ketika ia mampu menjadi
jembatan bagi 4.5 miliar orang yang selama ini terabaikan. Fokus saya
jelas: teknologi harus bekerja untuk menutup celah akses, bukan sekadar
menjadi tren inovasi tanpa dampak nyata pada kemanusiaan.

\bookmarksetup{startatroot}

\chapter{UAS-2 My Opinions}\label{uas-2-my-opinions}

\section{Transformasi Pendidikan Teknik di Era Kecerdasan Buatan: Dari
Hafalan ke Penciptaan
Nilai}\label{transformasi-pendidikan-teknik-di-era-kecerdasan-buatan-dari-hafalan-ke-penciptaan-nilai}

\subsection{Krisis Model Pendidikan
Konvensional}\label{krisis-model-pendidikan-konvensional}

Pembelajaran ilmu rekayasa saat ini sedang menghadapi pergeseran
paradigma yang sangat drastis. Menurut saya, model pembelajaran
konvensional yang hanya menekankan pada penguasaan sintaks atau hafalan
definisi sudah tidak lagi relevan. Problem utamanya adalah AI sekarang
sudah sangat mahir dalam melakukan tugas-tugas logika formal dan
komputasi yang jauh lebih cepat dari manusia.

Jika ujian mahasiswa masih berkisar pada menjawab soal-soal yang bisa
diselesaikan oleh Large Language Models dalam hitungan detik, maka nilai
dari pendidikan itu sendiri akan hilang. Mahasiswa bisa saja menggunakan
AI sebagai alat bantu tanpa benar-benar memahami logika di baliknya.
Yang menarik di sini adalah bagaimana kita seharusnya bergeser dari
sekadar menyelesaikan masalah teknis menuju kemampuan untuk memberikan
nilai nyata bagi masyarakat.

\subsection{Menggeser Fokus dari Bagaimana ke
Mengapa}\label{menggeser-fokus-dari-bagaimana-ke-mengapa}

Dalam konteks masalah kesehatan global yang saya bahas sebelumnya,
pendidikan teknik tidak boleh hanya mengajarkan cara membuat kode yang
efisien. Kita butuh lebih dari itu. Engineering bukan cuma soal teknis,
tapi soal empati. Seringkali kita terlalu fokus pada teknis koding tapi
lupa pada konteks manusia yang akan menggunakan sistem tersebut.

Pendidikan masa depan harus melatih mahasiswa untuk menjadi sosok yang
memiliki otoritas penuh atas teknologi yang mereka buat. AI seharusnya
menjadi pengganda kekuatan, bukan pengganti pemikiran kritis. Menurut
saya, kurikulum harus mulai mengintegrasikan pemahaman tentang dampak
sosial dan ekonomi dari sebuah sistem informasi, sehingga kita tidak
hanya membangun aplikasi yang keren secara teknis, tapi juga solutif
secara kemanusiaan.

\subsection{Ruang Kelas sebagai Miniatur Lingkungan
Profesional}\label{ruang-kelas-sebagai-miniatur-lingkungan-profesional}

Saya berpendapat bahwa ruang kelas harus bertransformasi menjadi
lingkungan kerja yang nyata. Mahasiswa tidak boleh lagi hanya menjadi
penerima materi yang pasif. Kita perlu didorong untuk memproduksi
artefak pengetahuan yang memiliki nilai nyata, bukan sekadar mengerjakan
tugas untuk mendapatkan nilai di atas kertas.

\begin{itemize}
\item
  Penguasaan materi sebaiknya divalidasi melalui proyek nyata yang bisa
  dirasakan manfaatnya oleh orang lain, bukan melalui ujian pilihan
  ganda.
\item
  Fokus penilaian harus dialihkan pada proses berpikir, manajemen stres,
  dan kemampuan adaptasi terhadap teknologi baru.
\item
  Trade-off yang sering muncul adalah antara kecepatan belajar dan
  kedalaman pemahaman. AI bisa mempercepat proses, tapi kedalaman logika
  tetap harus dibangun melalui refleksi yang jujur.
\item
  Menurut saya, portofolio yang berisi jurnal reflektif dan bukti
  penyelesaian masalah jauh lebih berharga daripada transkrip nilai yang
  hanya berisi angka-angka tanpa cerita di baliknya.
\end{itemize}

\subsection{Kesimpulan: Menjadi Engineer yang
Thoughtful}\label{kesimpulan-menjadi-engineer-yang-thoughtful}

Pada akhirnya, tujuan akhir dari pendidikan teknik di era AI adalah
membentuk manusia yang bijaksana dalam menggunakan alatnya. AI adalah
instrumen yang sangat kuat, tapi tanpa arah yang jelas dari manusianya,
teknologi tersebut bisa menjadi sia-sia atau bahkan merugikan.

Bagi saya, menjadi seorang rekayasawan sistem informasi berarti memikul
tanggung jawab moral untuk memastikan bahwa setiap baris kode yang kita
tulis memiliki dampak positif. Kita harus berani melangkah keluar dari
zona nyaman teknis dan mulai berpikir secara sistemik untuk
menyelesaikan tantangan global seperti kemiskinan dan krisis kesehatan.
Pendidikan yang berhasil adalah pendidikan yang mampu membuat
mahasiswanya merasa bahwa mereka punya peran penting dalam menulis
narasi masa depan umat manusia.

\bookmarksetup{startatroot}

\chapter{UAS-3 My Innovations}\label{uas-3-my-innovations}

\section{Arsitektur Sistem Informasi Kesehatan Terdistribusi: Solusi
Resilien untuk Skala
Global}\label{arsitektur-sistem-informasi-kesehatan-terdistribusi-solusi-resilien-untuk-skala-global}

Inovasi yang saya tawarkan berfokus pada pembangunan sistem backend yang
mampu menangani masalah kesehatan global di daerah dengan infrastruktur
terbatas. Mengingat sekitar 50\% penduduk bumi tidak memiliki akses
layanan kesehatan dasar, sistem ini dirancang dengan prinsip
desentralisasi untuk menjamin ketersediaan data di titik pelayanan
paling ujung.

\subsection{1. Optimasi Edge-AI dan Model
Quantization}\label{optimasi-edge-ai-dan-model-quantization}

Masalah utama di lapangan adalah ketergantungan pada koneksi cloud yang
seringkali tidak stabil. Untuk mengatasi ini, saya mengusulkan
penggunaan Edge-AI yang berjalan langsung di perangkat tenaga medis.

\begin{itemize}
\item
  Model AI menggunakan teknik quantization dan pruning agar bisa
  berjalan pada hardware dengan resource rendah tanpa memakan banyak
  daya atau memori.
\item
  Algoritma CNN dioptimalkan untuk melakukan klasifikasi citra medis
  sederhana atau deteksi awal gejala penyakit menular secara lokal.
\item
  Trade-off yang muncul adalah sedikit penurunan akurasi demi latensi
  yang lebih rendah. Menurut saya, lebih baik memiliki sistem yang
  memberikan hasil 85\% akurat secara instan di pelosok daripada 99\%
  akurat tapi butuh koneksi internet yang tidak pernah ada.
\item
  Data hasil deteksi disimpan dalam basis data lokal seperti SQLite dan
  hanya akan melakukan sinkronisasi ke server pusat saat mendeteksi
  adanya koneksi internet melalui protokol yang ringan.
\end{itemize}

\subsection{2. Event-Driven Architecture untuk Logistik
Medis}\label{event-driven-architecture-untuk-logistik-medis}

Kematian balita akibat diare yang mencapai 297 ribu jiwa per tahun
seringkali disebabkan oleh keterlambatan distribusi logistik. Saya
mengusulkan arsitektur berbasis event untuk mengotomatisasi rantai
pasok.

\begin{itemize}
\item
  Menggunakan pesan asinkron untuk memastikan koordinasi stok antara
  klinik primer dan gudang pusat tetap berjalan meskipun ada gangguan
  jaringan.
\item
  Sistem menggunakan algoritma prediktif berbasis time-series untuk
  menghitung reorder point secara otomatis berdasarkan tren penyakit
  musiman dan data historis di wilayah tersebut.
\item
  Pemanfaatan API Gateway dengan strategi caching yang agresif untuk
  meminimalkan beban server saat terjadi lonjakan permintaan data di
  daerah rawan penyakit.
\item
  Dari sisi implementasi, ini soal bagaimana sistem bisa melakukan
  ``graceful degradation'', yaitu tetap berfungsi pada level minimum
  saat komponen lain sedang bermasalah.
\end{itemize}

\subsection{3. Integrasi Skor Risiko Aktuaria dan AI
Finansial}\label{integrasi-skor-risiko-aktuaria-dan-ai-finansial}

Terdapat 2 miliar orang yang menghadapi beban finansial berat akibat
biaya kesehatan. Inovasi di sini adalah menggabungkan logika aktuaria
dengan machine learning untuk sistem proteksi sosial.

\begin{itemize}
\item
  Sistem melakukan scoring otomatis terhadap data ekonomi pasien untuk
  mengidentifikasi potensi biaya katastrofik sebelum pasien jatuh ke
  jurang kemiskinan.
\item
  Data tersebut dienkripsi dengan standar tinggi untuk menjaga privasi
  medis, namun tetap bisa diakses oleh pihak penyedia subsidi melalui
  mekanisme otorisasi yang ketat.
\item
  Yang menarik bagi saya adalah penggunaan analitik ini untuk
  menyeimbangkan beban subsidi pemerintah secara real-time, sehingga
  bantuan tidak salah sasaran.
\item
  Secara teknis, ini melibatkan pemrosesan batch data besar untuk
  melihat korelasi antara riwayat kesehatan dan stabilitas ekonomi rumah
  tangga di tingkat populasi.
\end{itemize}

\subsection{4. Strategi Deployment dan Keamanan
Data}\label{strategi-deployment-dan-keamanan-data}

Implementasi sistem ini tidak bisa dilakukan secara sekaligus. Saya
mengusulkan pendekatan bertahap yang fokus pada stabilitas sistem di
lingkungan yang keras.

\begin{itemize}
\item
  Tahap awal fokus pada pembuatan RESTful API yang efisien dengan
  overhead kecil agar hemat bandwidth.
\item
  Implementasi sistem keamanan berbasis token yang punya masa berlaku
  fleksibel sesuai dengan kondisi lapangan dan profil risiko akses di
  daerah konflik.
\item
  Penggunaan containerization seperti Docker untuk memudahkan deployment
  di berbagai jenis server lokal yang tersedia di puskesmas atau klinik
  terpencil.
\item
  Bagi saya, inovasi sebenarnya bukan pada seberapa canggih bahasa
  pemrograman yang dipakai, tapi pada seberapa tahan sistem tersebut
  saat menghadapi kegagalan infrastruktur nyata di dunia yang sedang
  krisis.
\end{itemize}

\bookmarksetup{startatroot}

\chapter{UAS-4 My Knowledge}\label{uas-4-my-knowledge}

\section{Landasan Teori dan Analisis Data dalam Rekayasa Kesehatan
Global}\label{landasan-teori-dan-analisis-data-dalam-rekayasa-kesehatan-global}

Untuk membangun sistem informasi yang berdampak pada masalah kesehatan
global, kita perlu memahami teori dasar di balik distribusi layanan dan
manajemen risiko ekonomi. Bagian ini merangkum dasar pengetahuan yang
menjadi fondasi inovasi saya dalam mengatasi krisis kesehatan bagi 4.5
miliar penduduk bumi.

\subsection{1. Teori Universal Health Coverage dan Gap Layanan
Dasar}\label{teori-universal-health-coverage-dan-gap-layanan-dasar}

Universal Health Coverage (UHC) adalah indikator utama untuk mengukur
apakah sebuah populasi memiliki akses ke layanan kesehatan esensial
tanpa mengalami kesulitan finansial. Masalahnya, laporan terbaru
menunjukkan bahwa kemajuan UHC telah mengalami stagnasi sejak tahun
2015.

\begin{itemize}
\item
  Saat ini, 4.5 miliar orang atau sekitar 50\% penduduk bumi masih belum
  terlayani oleh layanan kesehatan dasar secara penuh.
\item
  Kesenjangan ini paling terasa di negara berpenghasilan rendah, di mana
  satu dari dua orang kesulitan mengakses perawatan dasar seperti
  imunisasi atau pelayanan maternal.
\item
  Menurut saya, kegagalan pencapaian target SDG 3.8 bukan hanya karena
  kurangnya dana, tapi karena ketidakmampuan sistem informasi dalam
  memetakan kebutuhan di level akar rumput secara akurat.
\item
  Yang menarik untuk diperhatikan adalah bagaimana data statis
  seringkali gagal menangkap dinamika kesehatan di daerah konflik atau
  pelosok yang sangat membutuhkan intervensi cepat.
\end{itemize}

\subsection{2. Komputasi Tepi dan Optimasi Model AI untuk Skala
Global}\label{komputasi-tepi-dan-optimasi-model-ai-untuk-skala-global}

Secara teknis, menerapkan AI di daerah terpencil membutuhkan pemahaman
mendalam tentang komputasi tepi atau edge computing. Kita tidak bisa
memaksakan arsitektur cloud-heavy di wilayah yang infrastruktur
listriknya saja belum stabil.

\begin{itemize}
\item
  TinyML merupakan cabang dari machine learning yang memungkinkan
  eksekusi model pada perangkat dengan daya rendah dan memori terbatas
  melalui teknik kompresi model.
\item
  Implementasi di lapangan membutuhkan pemahaman tentang optimasi
  latensi agar sistem tetap responsif meski hardware yang digunakan
  sangat minimalis.
\item
  Trade-off utama dalam penerapan ini adalah antara kompleksitas model
  dan daya tahan baterai perangkat. Sebagai engineer, saya harus mampu
  memilih titik keseimbangan agar alat diagnostik tetap bisa dipakai
  seharian di lapangan tanpa perlu sering diisi daya.
\item
  Bagi saya, efisiensi algoritma di sini bukan lagi soal gaya-gayaan
  koding, tapi soal memastikan alat tersebut tetap menyala saat sedang
  dibutuhkan untuk menyelamatkan nyawa.
\end{itemize}

\subsection{3. Ekonomi Kesehatan dan Manajemen Risiko
Finansial}\label{ekonomi-kesehatan-dan-manajemen-risiko-finansial}

Kesehatan tidak bisa dipisahkan dari stabilitas ekonomi. Sekitar 2
miliar orang di dunia saat ini menghadapi beban finansial yang berat
karena biaya pengobatan, dan ratusan juta di antaranya terdorong ke
jurang kemiskinan setiap tahun.

\begin{itemize}
\item
  Masalah kesehatan global diperburuk oleh biaya tak terduga yang harus
  dibayar sendiri oleh pasien dari kantong mereka.
\item
  Manajemen risiko dalam sistem informasi kesehatan melibatkan
  penggunaan analitik prediktif untuk menghitung skor risiko ekonomi
  pasien berdasarkan data sosio-demografis.
\item
  Sebagai mahasiswa yang mempelajari aktuaria, saya melihat bahwa kita
  bisa membangun model proteksi sosial yang lebih cerdas dengan
  mengidentifikasi kelompok rentan sebelum krisis kesehatan benar-benar
  menghancurkan ekonomi mereka.
\item
  Problem utama yang sering saya temui adalah kebijakan subsidi yang
  seringkali bersifat satu untuk semua, padahal risiko setiap individu
  di wilayah yang berbeda memiliki profil yang sangat kontras.
\end{itemize}

\subsection{4. Analisis Data Penyakit yang Dapat
Dicegah}\label{analisis-data-penyakit-yang-dapat-dicegah}

Fokus pada kesehatan primer berarti fokus pada penyakit yang seharusnya
bisa dicegah melalui intervensi logistik dan edukasi yang tepat. Data
kematian balita menjadi parameter yang sangat penting di sini.

\begin{itemize}
\item
  Lebih dari 297 ribu anak balita meninggal setiap tahun akibat diare
  yang disebabkan oleh buruknya sanitasi dan air minum yang
  terkontaminasi.
\item
  Penanganan masalah ini membutuhkan sinkronisasi data real-time antara
  sistem informasi air bersih dan sistem kesehatan untuk memprediksi
  potensi wabah di sebuah wilayah.
\item
  Menurut saya, inovasi yang paling mendesak adalah membangun
  keterkaitan data antar berbagai sektor layanan dasar agar penanganan
  kesehatan tidak berjalan sendirian dalam silo.
\item
  Pada akhirnya, pengetahuan teknis tentang pengelolaan data besar hanya
  akan berguna jika kita tahu variabel mana yang paling krusial untuk
  dipantau demi mengurangi angka kematian prematur secara signifikan.
\end{itemize}

\bookmarksetup{startatroot}

\chapter{UAS-5 My Professional
Reviews}\label{uas-5-my-professional-reviews}

Untuk melAkukan review, seperti pada
\href{../My_Personal_Reviews/Doc.5.Mengevaluasi-Esai-Berdasarkan-Rubrik.pdf}{pendekatan
AI}, kita membutuhkan rubrik

\bookmarksetup{startatroot}

\chapter{Summary}\label{summary}

In summary, this book has no content whatsoever.

\bookmarksetup{startatroot}

\chapter*{References}\label{references}
\addcontentsline{toc}{chapter}{References}

\markboth{References}{References}

\phantomsection\label{refs}




\end{document}
